\subsection{Struttura della funzione e profilo energetico}

In Serverledge, ogni funzione è rappresentata dalla struttura 
\texttt{Function}, estesa in questo lavoro con i campi necessari al 
supporto delle varianti energetiche. I campi rilevanti per il sistema 
proposto sono i seguenti:

\begin{itemize}
    \item \texttt{LogicalName}: identificatore condiviso da tutte le 
    varianti appartenenti alla stessa funzione logica. Costituisce 
    il riferimento con cui il client invoca la funzione, 
    indipendentemente dalla variante che verrà selezionata a runtime.

    \item \texttt{VariantID}: identificatore univoco della specifica 
    variante all'interno del gruppo logico (ad esempio \texttt{base}, 
    \texttt{light}, \texttt{c}).

    \item \texttt{EnergyProfile}: struttura che mantiene la stima 
    energetica della variante, composta da due campi:
    \begin{itemize}
        \item \texttt{ColdStartJoule}: energia stimata per l'avvio a 
        freddo del container, comprensiva dell'inizializzazione del 
        runtime;
        \item \texttt{InvocationJoule}: energia media stimata per una 
        singola invocazione della funzione
    \end{itemize}
    \item \texttt{AllowApprox}: flag booleano che indica se la funzione 
    logica ammette l'utilizzo di varianti approssimate. La sua 
    attivazione è condizione necessaria affinché lo scheduler 
    energy-aware possa considerare varianti con accuratezza ridotta.
\end{itemize}

\subsubsection{Determinazione dei profili energetici iniziali}

I valori iniziali di \texttt{EnergyProfile} sono stati determinati 
offline tramite una campagna di misurazioni preliminari condotta con 
Kepler. La metodologia adottata tiene conto della natura delle metriche 
esposte da Kepler: trattandosi di contatori cumulativi, l'energia 
osservata al termine di $n$ invocazioni consecutive sullo stesso 
container rappresenta la somma dei contributi di tutte le esecuzioni, 
inclusa quella iniziale soggetta al cold start.

Per separare i due contributi energetici si è proceduto nel modo 
seguente. Dato un container su cui vengono eseguite $n$ invocazioni 
consecutive, si indica con $E_{tot}$ l'energia cumulativa totale 
rilevata da Kepler al termine dell'ultima invocazione e con $E_1$ 
l'energia della prima invocazione, che comprende sia il costo di 
inizializzazione del container che quello di esecuzione della funzione. 
Le invocazioni successive, eseguite a container già attivo, sono invece 
soggette al solo costo di esecuzione. Il costo medio di invocazione in 
warm start è quindi stimato come:

\begin{equation}
    E_{inv} = \frac{E_{tot} - E_1}{n - 1}
    \label{eq:inv-joule}
\end{equation}

e il costo di cold start come:

\begin{equation}
    E_{cold} = E_1 - E_{inv}
    \label{eq:cold-joule}
\end{equation}

La ripetizione delle misurazioni su un numero statisticamente 
significativo di invocazioni consente di attenuare la varianza 
introdotta dal sistema operativo e dal runtime, ottenendo stime 
robuste per entrambi i contributi. I valori così ottenuti vengono 
codificati nel file JSON di definizione delle varianti e costituiscono 
la stima iniziale del profilo energetico. Il \autoref{lst:fibonacci-json} 
mostra un esempio concreto per la funzione \texttt{fibonacci}, che 
espone tre varianti con profili energetici distinti.

\begin{lstlisting}[
    language=json,
    caption={File di definizione delle varianti per la funzione 
    \texttt{fibonacci}. Ciascuna variante specifica runtime, entry 
    point, profilo energetico iniziale e modello di output.},
    label={lst:fibonacci-json}
]
{
  "variants": [
    {
      "id": "base",
      "runtime": "python310",
      "entry_point": "fibonacci.handler",
      "src": "fibonacci.py",
      "energy": {
        "cold_start_joule": 0.06606,
        "invocation_joule": 0.000731579
      },
      "output": { "type": "error", "error_estimate": 0.0 },
      "is_approximate": false
    },
    {
      "id": "optimization-py",
      "runtime": "python310",
      "entry_point": "fibonacci_o.handler",
      "src": "fibonacci_o.py",
      "energy": {
        "cold_start_joule": 0.05640,
        "invocation_joule": 0.000494737
      },
      "output": { "type": "error", "error_estimate": 0.0 },
      "is_approximate": false
    },
    {
      "id": "c",
      "runtime": "native",
      "entry_point": "fibonacci",
      "src": "fibonacci",
      "energy": {
        "cold_start_joule": 0.008268422,
        "invocation_joule": 0.000171579
      },
      "output": { "type": "error", "error_estimate": 0.0 },
      "is_approximate": false
    }
  ]
}
\end{lstlisting}

I valori iniziali del profilo energetico sono progressivamente 
raffinati a runtime dal worker periodico, che aggiorna entrambi i valori sulla base delle misurazioni effettive 
raccolte durante l'esecuzione in produzione. I risultati numerici 
delle misurazioni preliminari per tutte le varianti sono riportati 
nel Capitolo~\ref{cap:RisultatiSperimentali}.